\documentclass[../../../include/open-logic-section]{subfiles}

\begin{document}

\olfileid{sth}{cardinals}{cp}
\olsection{Prinsip Cantor}

Élinga manèh \olref[ordinals][vn]{sec}. Nalika iku kita ngrembug
himpunan terurut becik lan ngusulaké supaya ana objek kang bisa
dadi wakilé pengurutan becik. Kanthi pamikiran iki, kita ngenalaké
ordinal, banjur ing
\olref[ordinals][ordtype]{ordtypesworklikeyouwant} kita nuduhaké
yèn ordinal tumindak kaya kang dikarepake, yaiku:
\[
	\ordtype{A, <} = \ordtype{B, \lessdot}
	\text{ yen lan mung yen } \tuple{A, <} \isomorphic \tuple{B, \lessdot}.
\]
Élinga luwih adoh manèh marang \olref[sfr][siz][equ]{sec}.
Ing kono, kanthi cara naif, kita ngenalaké pangertèn “ukuran”
himpunan. Tepaté, kita kandha yèn rong himpunan padha cacahé,
$\cardeq{A}{B}$, yèn lan mung yèn ana !!a{bijection} $f \colon A \to
B$. Pangertèn iki sacara intrinsik {luwih prasaja} tinimbang
pengurutan becik: kita mung migatèkaké !!{bijection}s, dudu
(kaya isomorfisme urutan) apa !!{bijection}s iku “njaga sawijining
struktur”.

Saka kéné tuwuh gagasan kang lumrah. Kaya kita ngenalaké objek
tartamtu, \emph{ordinal}, kanggo ngukur jinis pengurutan becik,
kita bisa ngenalaké objek tartamtu, \emph{kardinal}, kanggo ngukur
ukuran. Iku ancas bab iki.

Sadurungé njlentrehaké apa kardinal iku, kita perlu netepaké
prinsip kang kudu dituruti. Yèn $\card{X}$ nyathet kardinalitas
himpunan $X$, kita péngin supaya:
\[
	\card{A} = \card{B} \text{ yen lan mung yen } \cardeq{A}{B}.
\]
Iki bakal kita sebut \emph{Prinsip Cantor}, amarga Cantor mbokmenawa
wong kapisan kang nduwèni gagasan iki kanthi cetha. (Bab gegayutané
karo \emph{Prinsip Hume} bakal kita rembug luwih jembar ing
\olref[hp]{sec}.) Mula ancas kita yaiku netepaké $\card{X}$ kanggo
saben $X$ saéngga Prinsip Cantor bisa kawujud.

\end{document}
